% English draft based on sections/13_conclusion_ja.tex
\section{Conclusion: What Remains Open}

This paper has described the dynamics of intelligence, not prescribed its implementation. Its claims concern how inference, logs, stopping layers, thresholds, and external reference fields move.

\subsection{Three Logical Lines}

\textbf{First: non-closure and external stopping.} The inference operator $F$ has no internal fixed point. Therefore $L_F$ alone cannot close. Stopping must be externalized.

\textbf{Second: threshold mobility.} The distinction between responsive and silent worlds is threshold-dependent. $\mu_\theta$ is the capacity to move the threshold and reorganize the movement region.

\textbf{Third: three-layer dynamics.} The motion of $L_F$ leaves traces, producing $L_R$. The recursive motion of $L_R$ under $H_{ij}$ gives rise to $L_E$. The cooperation of $L_F$, $L_R$, and $L_E$ sustains intelligence.

\subsection{VED Principle in the Intelligence Layer}

VED describes reality in terms of motion arising from difference, not as static categories. This paper applies that principle to intelligence. Intelligence is not an object, capacity, or final state. It is a maintained dynamical state of inference, logs, and emergent stopping.

\begin{figure}[H]
\centering
\includegraphics[width=\linewidth]{fig9_cross_domain_homology.pdf}
\caption{Cross-domain homology. The same pattern of incomplete closure, emergent attractor, and sustainable dynamics appears across physics, life, consciousness, intelligence, and society.}
\label{fig:part2-cross-domain-homology}
\figurenote{The central homology kernel consists of incomplete closure, emergent attractor, and sustainable dynamics. Each domain card is a different implementation of this kernel.}{The figure does not erase domain differences. It identifies a repeated generative grammar while preserving the autonomy of each layer. Intelligence is special, but not isolated from the VED tree.}
\end{figure}

\subsection{Open Questions}

Several questions are intentionally left open. The meaning of understanding is passed to the philosophical layer. AI implementation and governance are outside this paper's descriptive scope. Quantifying the developmental degree of $L_E$ remains a future task. Formalizing a stopping functional analogous to closure energy is also left for later work.

\subsection{Final Statement}

Intelligence does not complete. It does not close. It is not optimized into a final state. Yet it persists. It persists because stopping is supplied externally and emergently, and because such stopping remains revisable.

Sapiens-like intelligence adds endogenous $\mu_\theta$: the ability to move the region in which stopping can be installed. It can err because it can revise; it can revise because it can move; it can persist because it does not close absolutely.

Theory ends here. The questions do not.

\bigskip
